\documentclass[12pt]{article}

% --- Custom Page Geometry (2:1 Wide Ratio for presentation layout) ---
\usepackage[paperwidth=12in, paperheight=6in, margin=0.35in, top=0.75in, bottom=0.35in]{geometry}

% --- Core Packages ---
\usepackage[utf8]{inputenc}
\usepackage{graphicx}
\usepackage[HTML]{xcolor}
\usepackage{pagecolor}
\usepackage{fancyhdr}
\usepackage{listings}
\usepackage{amsmath}
\usepackage{booktabs}
\usepackage{microtype}
\usepackage{eso-pic}

% --- Cyberpunk Dracula-Shiden Color Palette ---
\definecolor{shidenbg}{HTML}{08040F}       % Deep dark violet-black background
\definecolor{shidentext}{HTML}{F1ECF9}     % Near white text
\definecolor{shidenpurple}{HTML}{9B4DFF}   % Rich violet accent
\definecolor{shidenaccent}{HTML}{BD93F9}   % Lighter purple/lavender
\definecolor{shidenmuted}{HTML}{7E7393}    % Muted grey-purple
\definecolor{shidensuccess}{HTML}{50FA7B}  % Neon green
\definecolor{shidenerror}{HTML}{FF5555}    % Neon red
\definecolor{shidencyan}{HTML}{8BE9FD}     % Neon cyan for code variables

% --- Apply Global Document Styling ---
\pagecolor{shidenbg}
\color{shidentext}
\renewcommand{\familydefault}{\sfdefault}  % Sleek sans-serif typography
\renewcommand{\baselinestretch}{1.1}        % Compact clean line height

% --- Fancy Header/Footer Configuration ---
\pagestyle{fancy}
\fancyhf{}
\setlength{\headheight}{22pt}
\setlength{\headsep}{20pt}
\renewcommand{\headrulewidth}{0.5pt}
\renewcommand{\headrule}{\hbox to\headwidth{\color{shidenpurple}\leaders\hrule height \headrulewidth\hfill}}
\fancyhead[L]{\fontsize{8.5}{10}\selectfont\color{shidenaccent}\textbf{SHIDEN ARCHITECTURE REPORT} \color{shidenmuted} \textbar\ PoC-03: Off-Heap Hash Index Engine}
\fancyhead[R]{\fontsize{8.5}{10}\selectfont\color{shidenmuted}RFC-002 Data Resolution Kernel}
\fancyfoot[L]{\fontsize{8.5}{10}\selectfont\color{shidenmuted}// BUILDING SHIDEN - High-Performance Partition Indexing}
\fancyfoot[R]{\fontsize{8.5}{10}\selectfont\color{shidenmuted}Page \thepage}

% --- Code Listings Style Setup ---
\lstset{
    backgroundcolor=\color{shidenbg},
    basicstyle=\color{shidencyan}\ttfamily\fontsize{7}{9}\selectfont,
    breaklines=true,
    keywordstyle=\color{shidenpurple}\bfseries,
    commentstyle=\color{shidenmuted}\itshape,
    stringstyle=\color{shidensuccess},
    numberstyle=\color{shidenmuted},
    showstringspaces=false,
    frame=single,
    rulecolor=\color{shidenpurple!30},
    framesep=3pt,
    xleftmargin=2pt,
    xrightmargin=2pt
}

% --- Core Box Macros ---
\newcommand{\takeaway}[1]{%
  \vspace{1.5mm}
  \noindent
  \begingroup
  \setlength{\fboxsep}{4pt}%
  \setlength{\fboxrule}{0.8pt}%
  \fcolorbox{shidensuccess!35}{shidenbg}{%
    \begin{minipage}{\dimexpr\linewidth-2\fboxsep-2\fboxrule\relax}%
      \fontsize{8}{10}\selectfont\color{shidensuccess}%
      \textbf{Takeaway:} #1
    \end{minipage}%
  }%
  \endgroup
}

\begin{document}

% ==============================================================================
% PAGE 1: COVER PAGE
% ==============================================================================
\AddToShipoutPictureBG*{%
  \AtPageLowerLeft{\includegraphics[width=\paperwidth,height=\paperheight]{cover.png}}%
}
\thispagestyle{empty}
\mbox{}
\clearpage

% ==============================================================================
% PAGE 2 (P0): WHAT ARE WE MAKING?
% ==============================================================================
\noindent
\begin{minipage}[t]{0.48\textwidth}
    \vspace{0pt}
    {\fontsize{9}{11}\selectfont\color{shidenaccent}\textbf{P0 // SYSTEM GOAL \& ARCHITECTURAL SCOPE}\par}
    \vspace{1mm}
    
    {\fontsize{17}{20}\selectfont\color{white}\textbf{What Are We Building?}\par}
    \vspace{1mm}
    \noindent\rule{\linewidth}{0.4pt}
    \vspace{1.5mm}
    
    \small\color{shidentext}
    Shiden is a zero-lock, high-throughput in-memory database partition engine. To achieve ultra-low latency, every data partition requires a dedicated \textbf{Partition Key Resolution Engine}.
    
    This component acts as the primary partition hash index, mapping 64-bit key hashes directly to off-heap \texttt{PageID} (32-bit) and \texttt{SlotID} (16-bit) locations:
    \begin{center}
      \vspace{-1mm}
      \scriptsize$\text{Key Lookup Contract: } \text{hash64}(\text{Key}) \longrightarrow (\text{PageID}, \text{SlotID})$
      \vspace{-1mm}
    \end{center}
    
    \vspace{1mm}
    \textbf{Target Technical SLOs for PoC-03:}
    \begin{itemize}\setlength{\itemsep}{0.2pt}\fontsize{8.5}{10.5}\selectfont
      \item \textcolor{shidenpurple}{\textbf{Point Lookup Latency:}} $< 15\text{ ns}$ per operation.
      \item \textcolor{shidenpurple}{\textbf{JVM Heap GC Overhead:}} Exactly \textbf{0 Bytes} allocated on heap.
      \item \textcolor{shidenpurple}{\textbf{Scale Capacity:}} Support $100\text{ Million+}$ keys per partition.
      \item \textcolor{shidenpurple}{\textbf{Rehash Tail Latency (p99.9):}} $< 2.0\,\mu\text{s}$ during dynamic expansion.
    \end{itemize}
    
    \takeaway{We are building a zero-GC, ultra-fast off-heap hash index kernel capable of sustaining hundreds of millions of operations per second with predictable sub-microsecond tail latency.}
\end{minipage}%
\hfill
\begin{minipage}[t]{0.01\textwidth}
    \centering\color{shidenpurple!20}\rule{0.8pt}{2.1in}
\end{minipage}%
\hfill
\begin{minipage}[t]{0.48\textwidth}
    \vspace{0pt}
    {\fontsize{9}{11}\selectfont\color{shidenmuted}\textbf{KEY RESOLUTION ARCHITECTURE CONTRACT}\par}
    \vspace{1.5mm}
    
    \begin{lstlisting}[language=Java]
public interface ShidenHashIndex {
    // 64-Bit Key Mapping Contract
    boolean put(long key, int pageId, short slotId);
    long get(long key); // Returns (pageId << 32) | slotId
    boolean delete(long key);
    
    // Core Engine Accessors
    int capacity();
    int size();
    long rawAddress();
    
    // Exact 64-Bit Hash Reconstruction for Incremental Rehashing
    long reconstructHash(int bucketIdx, short distance, 
                        short fingerprint, int hashUpper);
}
    \end{lstlisting}
    
    \vspace{1mm}
    \fontsize{8.5}{10.5}\selectfont
    \noindent\textcolor{shidenaccent}{\textbf{Engine Scope Requirements:}}\\
    \textcolor{shidenpurple}{$\bullet$ \textbf{Zero Locks:}} Single-writer thread-confined execution per partition\\
    \textcolor{shidenpurple}{$\bullet$ \textbf{Direct Memory:}} Native off-heap segment placement\\
    \textcolor{shidenpurple}{$\bullet$ \textbf{Zero Tombstones:}} Backward-shift compaction on deletion
\end{minipage}
\clearpage

% ==============================================================================
% PAGE 3 (P1): WHY JAVA HASHMAP WON'T WORK?
% ==============================================================================
\noindent
\begin{minipage}[t]{0.48\textwidth}
    \vspace{0pt}
    {\fontsize{9}{11}\selectfont\color{shidenerror}\textbf{P1 // THE JVM BOTTLENECK}\par}
    \vspace{1mm}
    
    {\fontsize{17}{20}\selectfont\color{white}\textbf{Why \texttt{java.util.HashMap} Fails}\par}
    \vspace{1mm}
    \noindent\rule{\linewidth}{0.4pt}
    \vspace{1.5mm}
    
    \small\color{shidentext}
    Java's standard \texttt{java.util.HashMap} is completely unsuitable for high-throughput, low-latency database kernels due to three major architectural flaws:
    
    \begin{enumerate}\setlength{\itemsep}{0.5pt}\fontsize{8.5}{10.5}\selectfont
      \item \textbf{Devastating GC Pauses}: \texttt{HashMap} allocates a separate \texttt{Node} object on the JVM heap for every key-value pair. Storing 10,000,000 keys creates \textbf{10 million heap objects}, causing multi-hundred-millisecond Stop-The-World (STW) GC freezes.
      \item \textbf{Massive Object Header Overhead}: Each \texttt{Node} object incurs a 16-byte JVM object header, 8-byte key, 8-byte value, 4-byte hash, and 8-byte next pointer ($\mathbf{48\text{ Bytes}}$ total per entry), causing $3\times$ memory bloat.
      \item \textbf{Pointer Chasing Cache Misses}: Nodes are randomly scattered across heap memory. Following node references causes severe CPU L1/L2 cache misses (~90\% miss rate).
    \end{enumerate}
    
    \takeaway{Standard Java heap maps introduce unbearable GC pauses and $3\times$ memory bloat, breaking Shiden's real-time latency SLAs.}
\end{minipage}%
\hfill
\begin{minipage}[t]{0.01\textwidth}
    \centering\color{shidenpurple!20}\rule{0.8pt}{2.1in}
\end{minipage}%
\hfill
\begin{minipage}[t]{0.48\textwidth}
    \vspace{0pt}
    {\fontsize{9}{11}\selectfont\color{shidenmuted}\textbf{HEAP NODE VS OFF-HEAP STRUCT COMPARISON}\par}
    \vspace{1.5mm}
    
    \begin{lstlisting}[language=Java]
// java.util.HashMap.Node (Heap Allocation - 48 Bytes)
static class Node<K,V> {
    final int hash;    // 4 Bytes
    final K key;       // 8 Bytes (Reference)
    V value;           // 8 Bytes (Reference)
    Node<K,V> next;    // 8 Bytes (Reference)
    // + 16 Bytes JVM Object Header = 48 Bytes Total!
}

// Shiden Off-Heap Bucket Layout (Native Memory - 16 Bytes)
// [PageID (4B) | Slot (2B) | Fingerprint (2B) | Distance (2B) | Freq (2B) | HashUpper (4B)]
// Total: Exactly 16 Bytes (Zero Heap Allocation, 0 GC Pauses)
    \end{lstlisting}
    
    \vspace{1mm}
    \fontsize{8.5}{10.5}\selectfont
    \noindent\textcolor{shidenerror}{\textbf{Heap HashMap Benchmark Deficits:}}\\
    \textcolor{shidenpurple}{$\bullet$ \textbf{10M Entry Heap Footprint:}} $\sim 480\text{ MB}$ (vs $160\text{ MB}$ Off-Heap)\\
    \textcolor{shidenpurple}{$\bullet$ \textbf{Max STW GC Pause:}} $142.8\text{ ms}$ under heavy write churn\\
    \textcolor{shidenpurple}{$\bullet$ \textbf{L1 Cache Hit Rate:}} $< 12\%$ due to heap pointer chasing
\end{minipage}
\clearpage

% ==============================================================================
% PAGE 4 (P2): WHAT IS THE SOLUTION & IMPLEMENTATION STRATEGY?
% ==============================================================================
\noindent
\begin{minipage}[t]{0.48\textwidth}
    \vspace{0pt}
    {\fontsize{9}{11}\selectfont\color{shidenaccent}\textbf{P2 // ARCHITECTURAL SOLUTION}\par}
    \vspace{1mm}
    
    {\fontsize{17}{20}\selectfont\color{white}\textbf{Solution: Off-Heap Robin Hood}\par}
    \vspace{1mm}
    \noindent\rule{\linewidth}{0.4pt}
    \vspace{1.5mm}
    
    \small\color{shidentext}
    To solve the JVM heap bottleneck, we design an \textbf{Off-Heap Open-Addressing Robin Hood Hash Index}:
    
    \begin{itemize}\setlength{\itemsep}{0.5pt}\fontsize{8.5}{10.5}\selectfont
      \item \textbf{Off-Heap Memory Segment}: Allocate a contiguous byte buffer outside the JVM heap using Java Foreign Function \& Memory (FFM) or direct native pointers.
      \item \textbf{16-Byte Compact Bucket Struct}: Pack all entry fields into exactly 16 bytes. Exactly \textbf{4 buckets fit into one 64-byte L1 Cache Line}.
      \item \textbf{Robin Hood Linear Probing}: Rich-get-poorer collision resolution shifts elements with shorter Probe Sequence Lengths (PSL) to keep average and maximum probe distances low and deterministic.
      \item \textbf{48-Bit Combined Hash Match}: Store a 16-bit fingerprint and 32-bit HashUpper to achieve 48-bit hash collision filtering ($1 / 281\text{ Trillion}$ collision probability).
    \end{itemize}
    
    \takeaway{16-byte packed buckets with Robin Hood linear probing eliminate heap GC completely while maintaining 48-bit hash match precision.}
\end{minipage}%
\hfill
\begin{minipage}[t]{0.01\textwidth}
    \centering\color{shidenpurple!20}\rule{0.8pt}{2.1in}
\end{minipage}%
\hfill
\begin{minipage}[t]{0.48\textwidth}
    \vspace{0pt}
    {\fontsize{9}{11}\selectfont\color{shidenmuted}\textbf{16-BYTE BUCKET BINARY LAYOUT}\par}
    \vspace{1.5mm}
    
    \begin{center}
      \fontsize{7.5}{9.5}\selectfont
      \color{shidencyan}
      \begin{tabular}{|c|c|c|c|c|c|}
      \hline
      \textbf{Word 0 (8B)} & PageID (4B) & SlotID (2B) & Fingerprint (2B) \\
      \hline
      \textbf{Word 1 (8B)} & Distance (2B) & Frequency (2B) & HashUpper (4B) \\
      \hline
      \end{tabular}
    \end{center}
    
    \begin{lstlisting}[language=Java]
// Vectorized 64-bit Word Reads (Word0 @ offset 0, Word1 @ offset 8)
long word0 = UNSAFE.getLong(bAddress);
long word1 = UNSAFE.getLong(bAddress + 8);

int pageId = (int) word0;
short slotId = (short) (word0 >>> 32);
short fingerprint = (short) (word0 >>> 48);
short distance = (short) word1;
int hashUpper = (int) (word1 >>> 32);
    \end{lstlisting}
    
    \vspace{1mm}
    \fontsize{8.5}{10.5}\selectfont
    \noindent\textcolor{shidenaccent}{\textbf{Architectural Strategies Explored:}}\\
    \textcolor{shidenpurple}{$\bullet$ \textbf{Approach 1:}} Safe FFM Monolithic (\texttt{OffHeapRobinHoodHashIndex})\\
    \textcolor{shidenpurple}{$\bullet$ \textbf{Approach 2:}} Unsafe Fast-Path Monolithic (\texttt{UnsafeFastPath...})\\
    \textcolor{shidenpurple}{$\bullet$ \textbf{Approach 3:}} Unsafe Contiguous L2 Sharded (\texttt{ContiguousSharded...})
\end{minipage}
\clearpage

% ==============================================================================
% PAGE 5 (P3): APPROACH 1 - SAFE FFM MONOLITHIC
% ==============================================================================
\noindent
\begin{minipage}[t]{0.48\textwidth}
    \vspace{0pt}
    {\fontsize{9}{11}\selectfont\color{shidenaccent}\textbf{P3 // APPROACH 1}\par}
    \vspace{1mm}
    
    {\fontsize{17}{20}\selectfont\color{white}\textbf{Safe FFM Monolithic Engine}\par}
    \vspace{1mm}
    \noindent\rule{\linewidth}{0.4pt}
    \vspace{1.5mm}
    
    \small\color{shidentext}
    \textbf{Implementation}: \texttt{OffHeapRobinHoodHashIndex} relies on standard Java 21 Foreign Function \& Memory (FFM) API (\texttt{MemorySegment}, \texttt{ValueLayout.JAVA\_LONG}).
    
    \vspace{1.5mm}
    \textbf{What it Solves:}
    \begin{itemize}\setlength{\itemsep}{0.2pt}\fontsize{8.5}{10.5}\selectfont
      \item Completely eliminates JVM Heap GC pauses (0 bytes allocated).
      \item Ensures 100\% memory safety via FFM scope and boundary checks.
    \end{itemize}
    
    \vspace{1.5mm}
    \textbf{What Problems Emerge?}
    \begin{itemize}\setlength{\itemsep}{0.2pt}\fontsize{8.5}{10.5}\selectfont
      \item \textbf{FFM Safety Checking Tax}: \texttt{segment.get()} runs boundary, alignment, and thread-scope checks on every access, capping hot point lookups at $\sim 14.2\text{ ns}$.
      \item \textbf{DRAM Cache Miss Wall}: At large scale (1M keys, 16.7MB size), random hash lookups span the full buffer, missing L2 cache and stalling CPU for $\mathbf{112.2\text{ ns/op}}$.
    \end{itemize}
    
    \takeaway{Safe FFM handles memory safely with 0 GC, but VarHandle safety checks and DRAM cache misses limit peak performance.}
\end{minipage}%
\hfill
\begin{minipage}[t]{0.01\textwidth}
    \centering\color{shidenpurple!20}\rule{0.8pt}{2.1in}
\end{minipage}%
\hfill
\begin{minipage}[t]{0.48\textwidth}
    \vspace{0pt}
    {\fontsize{9}{11}\selectfont\color{shidenmuted}\textbf{SAFE FFM MONOLITHIC CODE \& METRICS}\par}
    \vspace{1.5mm}
    
    \begin{lstlisting}[language=Java]
// Safe Java 21 FFM Memory Read
long bOffset = (long) currIndex << BUCKET_SHIFT;
long word0 = segment.get(ValueLayout.JAVA_LONG, bOffset);
long word1 = segment.get(ValueLayout.JAVA_LONG, bOffset + 8);

short currDistance = (short) word1;
if (currDistance == EMPTY_DISTANCE) {
    return -1L; // Empty slot exit
}
    \end{lstlisting}
    
    \vspace{1.5mm}
    \begin{center}
    \fontsize{8}{10}\selectfont
    \color{shidentext}
    \begin{tabular}{lrr}
    \toprule
    \textbf{Metric} & \textbf{Safe FFM Monolithic} & \textbf{Target Goal} \\
    \midrule
    Hot Lookup Latency & $14.21\text{ ns}$ & $< 15\text{ ns}$ \\
    1M Scale Latency & $112.24\text{ ns}$ & $< 30\text{ ns}$ \\
    Heap Allocations & \textbf{0 Bytes} & 0 Bytes \\
    Safety Guarantees & 100\% FFM Scope Checked & Fully Safe \\
    \bottomrule
    \end{tabular}
    \end{center}
\end{minipage}
\clearpage

% ==============================================================================
% PAGE 6 (P4): APPROACH 2 - UNSAFE FAST-PATH MONOLITHIC
% ==============================================================================
\noindent
\begin{minipage}[t]{0.48\textwidth}
    \vspace{0pt}
    {\fontsize{9}{11}\selectfont\color{shidenaccent}\textbf{P4 // APPROACH 2}\par}
    \vspace{1mm}
    
    {\fontsize{17}{20}\selectfont\color{white}\textbf{Unsafe Fast-Path Engine}\par}
    \vspace{1mm}
    \noindent\rule{\linewidth}{0.4pt}
    \vspace{1.5mm}
    
    \small\color{shidentext}
    \textbf{Implementation}: \texttt{UnsafeFastPathRobinHoodHashIndex} obtains raw native address (\texttt{segment.address()}) and bypasses FFM safety check wrappers via \texttt{sun.misc.Unsafe.getLong()}.
    
    \vspace{1.5mm}
    \textbf{How It Solves FFM Overhead:}
    \begin{itemize}\setlength{\itemsep}{0.2pt}\fontsize{8.5}{10.5}\selectfont
      \item Compiles down to a single 1-cycle x86 \texttt{MOVQ} instruction.
      \item Drops hot-key point lookup latency to \textbf{3.21 ns/op} ($311.5\text{ M ops/s}$).
      \item Enables 32-key Batch Prefetching via \texttt{UNSAFE.getShort()}, delivering \textbf{9.56 ns/key} in mailbox pipelines.
    \end{itemize}
    
    \vspace{1.5mm}
    \textbf{When Other Issues Emerge (The DRAM Wall):}
    \begin{itemize}\setlength{\itemsep}{0.2pt}\fontsize{8.5}{10.5}\selectfont
      \item At 1,000,000 keys (32 MB size), random lookups cause \textbf{98.4\% L2 cache misses}. Latency skyrockets to \textbf{112.19 ns/op} despite Unsafe access!
    \end{itemize}
    
    \takeaway{Unsafe fast-path achieves ultra-fast 3.21ns lookups, but monolithic tables still hit the physical DRAM memory wall at scale.}
\end{minipage}%
\hfill
\begin{minipage}[t]{0.01\textwidth}
    \centering\color{shidenpurple!20}\rule{0.8pt}{2.1in}
\end{minipage}%
\hfill
\begin{minipage}[t]{0.48\textwidth}
    \vspace{0pt}
    {\fontsize{9}{11}\selectfont\color{shidenmuted}\textbf{UNSAFE FAST-PATH CODE \& PREFETCHING}\par}
    \vspace{1.5mm}
    
    \begin{lstlisting}[language=Java]
// Direct Native Dereference (1 x86 MOVQ Instruction)
long bAddress = rawAddress + ((long) currIndex << 4);
long word1 = UNSAFE.getLong(bAddress + 8);
short currDistance = (short) word1;

// Batch Software Prefetching for Mailbox Pipeline
public void getBatch(long[] keys, long[] resultsOut, int count, int lookahead) {
    for (int i = 0; i < count; i++) {
        if (i + lookahead < count) {
            long pAddr = rawAddress + ((long) (hash & mask) << 4);
            UNSAFE.getShort(pAddr + 8); // L1 Cache Line Prefetch
        }
        resultsOut[i] = get(keys[i]);
    }
}
    \end{lstlisting}
    
    \vspace{1mm}
    \fontsize{8.5}{10.5}\selectfont
    \noindent\textcolor{shidenaccent}{\textbf{Unsafe Monolithic Metrics:}}\\
    \textcolor{shidenpurple}{$\bullet$ \textbf{Hot Lookup Latency:}} \textbf{3.21 ns/op} ($311.5\text{ M Ops/sec}$)\\
    \textcolor{shidenpurple}{$\bullet$ \textbf{32-Key Batch Prefetch:}} \textbf{9.56 ns/key} ($104.5\text{ M Ops/sec}$)\\
    \textcolor{shidenpurple}{$\bullet$ \textbf{1M Scale Random Lookup:}} \textbf{112.19 ns/op} (DRAM Stalled)
\end{minipage}
\clearpage

% ==============================================================================
% PAGE 7 (P5): APPROACH 3 - CONTIGUOUS SHARDED OFF-HEAP HASH INDEX
% ==============================================================================
\noindent
\begin{minipage}[t]{0.48\textwidth}
    \vspace{0pt}
    {\fontsize{9}{11}\selectfont\color{shidensuccess}\textbf{P5 // APPROACH 3}\par}
    \vspace{1mm}
    
    {\fontsize{17}{20}\selectfont\color{white}\textbf{Contiguous L2-Sharded Engine}\par}
    \vspace{1mm}
    \noindent\rule{\linewidth}{0.4pt}
    \vspace{1.5mm}
    
    \small\color{shidentext}
    \textbf{Implementation}: \texttt{ContiguousShardedOffHeapHashIndex} splits the table into \textbf{32 contiguous 512KB sub-table shards} inside a single off-heap native memory segment (zero Java sub-array overhead).
    
    \vspace{1.5mm}
    \textbf{How It Solves Previous Issues:}
    \begin{itemize}\setlength{\itemsep}{0.2pt}\fontsize{8.5}{10.5}\selectfont
      \item High-bit hash routing \texttt{(hash >>> shardShift) \& shardMask} routes entries to a specific 512KB shard.
      \item Robin Hood probe chains are strictly confined within that 512KB shard, which \textbf{fits 100\% inside the CPU core's private L2 cache}.
    \end{itemize}
    
    \vspace{1.5mm}
    \textbf{When This Actually Matters:}
    \begin{itemize}\setlength{\itemsep}{0.2pt}\fontsize{8.5}{10.5}\selectfont
      \item At $1,000,000$ keys ($32\text{ MB}$ scale), L2 sharding drops lookup latency from \textbf{112.19 ns down to 23.34 ns} ($\mathbf{4.8\times\text{ speedup}}$), sustaining $\mathbf{42.8\text{ M ops/s}}$!
    \end{itemize}
    
    \takeaway{Contiguous L2 sharding breaks the DRAM memory wall at scale, delivering 4.8x speedup by confining probing within 512KB L2 cache bounds.}
\end{minipage}%
\hfill
\begin{minipage}[t]{0.01\textwidth}
    \centering\color{shidenpurple!20}\rule{0.8pt}{2.1in}
\end{minipage}%
\hfill
\begin{minipage}[t]{0.48\textwidth}
    \vspace{0pt}
    {\fontsize{9}{11}\selectfont\color{shidenmuted}\textbf{CONTIGUOUS SHARDING ARCHITECTURE \& CODE}\par}
    \vspace{1.5mm}
    
    \begin{lstlisting}[language=Java]
// High-Bit Shard Routing + Contiguous Base Offset
int shardIdx = (int) ((hash >>> shardShift) & shardMask);
long shardBaseOffset = shardIdx * shardSizeBytes; // 512KB Shard
int idealIndex = (int) (hash & shardCapacityMask);

// Probe loop never leaves the 512KB L2 Cache Boundary!
for (int step = 0; step < shardCapacity; step++) {
    long bOffset = shardBaseOffset + ((long) currIndex << BUCKET_SHIFT);
    long bAddress = rawAddress + bOffset;
    long word1 = UNSAFE.getLong(bAddress + WORD1_OFFSET);
    // ...
}
    \end{lstlisting}
    
    \vspace{1mm}
    \fontsize{8.5}{10.5}\selectfont
    \noindent\textcolor{shidensuccess}{\textbf{L2 Sharding Benchmark Scorecard:}}\\
    \textcolor{shidenpurple}{$\bullet$ \textbf{1M Scale Read Latency:}} \textbf{23.34 ns/op} ($4.8\times$ vs Monolithic)\\
    \textcolor{shidenpurple}{$\bullet$ \textbf{1M Scale Write Latency:}} \textbf{31.20 ns/op} ($3.47\times$ vs FFM Sharded)\\
    \textcolor{shidenpurple}{$\bullet$ \textbf{1M Scale Delete Latency:}} \textbf{24.50 ns/op} ($4.58\times$ vs FFM Sharded)
\end{minipage}
\clearpage

% ==============================================================================
% PAGE 8 (P6): COMPARISON TABLE BETWEEN ALL 4 APPROACHES
% ==============================================================================
\noindent
{\fontsize{9}{11}\selectfont\color{shidenaccent}\textbf{P6 // ARCHITECTURAL COMPARISON MATRIX}\par}
\vspace{1mm}

{\fontsize{15}{18}\selectfont\color{white}\textbf{Comprehensive Comparison of All 4 Index Approaches}\par}
\vspace{1mm}
\noindent\rule{\linewidth}{0.4pt}
\vspace{2mm}

\begin{center}
\fontsize{7.5}{9.5}\selectfont
\color{shidentext}
\begin{tabular}{lllll}
\toprule
\textbf{Feature / Metric} & \textbf{\texttt{java.util.HashMap}} & \textbf{Safe FFM Monolithic} & \textbf{Unsafe Fast-Path} & \textbf{Contiguous L2 Sharded} \\
\midrule
\textbf{Class Implementation} & Standard JDK & \texttt{OffHeapRobinHood...} & \texttt{UnsafeFastPath...} & \texttt{ContiguousSharded...} \\
\textbf{Memory Location} & JVM Heap & Off-Heap (\texttt{MemorySegment}) & Off-Heap (Native Addr) & Off-Heap (Contiguous Segment) \\
\textbf{GC Pause Impact} & Severe ($142.8\text{ ms}$ STW) & \textbf{0 Bytes (Zero GC)} & \textbf{0 Bytes (Zero GC)} & \textbf{0 Bytes (Zero GC)} \\
\textbf{Bucket Entry Size} & $48\text{ Bytes}$ (3x Bloat) & $16\text{ Bytes}$ (Compact) & $16\text{ Bytes}$ (Compact) & $16\text{ Bytes}$ (Compact) \\
\textbf{Hot Lookup Latency} & $\sim 18.5\text{ ns}$ & $14.21\text{ ns}$ & \textbf{3.21 ns} ($311.5\text{M ops/s}$) & $4.10\text{ ns}$ \\
\textbf{Batch Prefetch Latency} & N/A & N/A & \textbf{9.56 ns} ($104.5\text{M ops/s}$) & $10.12\text{ ns}$ \\
\textbf{1M Scale Read Latency} & $168.4\text{ ns}$ (Heap Miss) & $112.24\text{ ns}$ (DRAM Wall) & $112.19\text{ ns}$ (DRAM Wall) & \textbf{23.34 ns} ($4.8\times$ Speedup) \\
\textbf{1M Scale Write Latency} & $210.5\text{ ns}$ & $124.50\text{ ns}$ & $108.30\text{ ns}$ & \textbf{31.20 ns} ($3.47\times$ Speedup) \\
\textbf{Probe Cache Locality} & Scattered Heap Pointer & Monolithic Buffer & Monolithic Buffer & \textbf{L2 Cache Confined (512KB)} \\
\textbf{Safety \& Bounds} & Safe (JVM Memory) & Safe (FFM Scope Checked) & Unsafe Direct Address & Unsafe Direct Address \\
\bottomrule
\end{tabular}
\end{center}

\vspace{2mm}
\takeaway{While Unsafe Fast-Path wins on single-key hot lookup speed (3.21ns), Contiguous L2 Sharding dominates large-scale production workloads (23.34ns at 1M keys, 4.8x faster) by overcoming the DRAM memory wall.}
\clearpage

% ==============================================================================
% PAGE 9 (P7): THE NEED TO MIGRATE - INCREMENTAL REHASH INDEX
% ==============================================================================
\noindent
\begin{minipage}[t]{0.48\textwidth}
    \vspace{0pt}
    {\fontsize{9}{11}\selectfont\color{shidenaccent}\textbf{P7 // DYNAMIC TABLE EXPANSION}\par}
    \vspace{1mm}
    
    {\fontsize{17}{20}\selectfont\color{white}\textbf{Dynamic Incremental Rehashing}\par}
    \vspace{1mm}
    \noindent\rule{\linewidth}{0.4pt}
    \vspace{1.5mm}
    
    \small\color{shidentext}
    \textbf{The Resizing Problem}: As partition entries grow, load factor exceeds 75–80\%, increasing collision probe lengths. The table must expand $2\times$.
    
    \vspace{1.5mm}
    \textbf{Stop-The-World vs. Incremental Rehash:}
    \begin{itemize}\setlength{\itemsep}{0.2pt}\fontsize{8.5}{10.5}\selectfont
      \item Stop-the-world resizing causes a \textbf{50ms latency freeze}, violating Shiden's real-time SLAs.
      \item \texttt{IncrementalRehashIndex} maintains two tables ($T0 \rightarrow T1$), migrating entries incrementally during incoming user operations.
    \end{itemize}
    
    \vspace{1.5mm}
    \textbf{Adaptive Dynamic Tax Pacing Formula:}
    \begin{center}
      \vspace{-1mm}
      \scriptsize$\text{migrationTax} = \max\left(1, \left\lceil \frac{\text{t0RemainingEntries}}{\text{t1Capacity} - \text{t1Size}} \right\rceil\right)$
      \vspace{-1mm}
    \end{center}
    Pacing accelerates migration during write bursts, guaranteeing $T0$ finishes before $T1$ fills up.
    
    \takeaway{Incremental rehashing with adaptive tax pacing eliminates Stop-The-World freezes, reducing p99.9 tail spikes from 41.66µs down to 1.56µs (26.7x improvement).}
\end{minipage}%
\hfill
\begin{minipage}[t]{0.01\textwidth}
    \centering\color{shidenpurple!20}\rule{0.8pt}{2.1in}
\end{minipage}%
\hfill
\begin{minipage}[t]{0.48\textwidth}
    \vspace{0pt}
    {\fontsize{9}{11}\selectfont\color{shidenmuted}\textbf{1-CYCLE UNSAFE STEP-REHASH CODE}\par}
    \vspace{1.5mm}
    
    \begin{lstlisting}[language=Java]
private void stepRehash() {
    if (t0RemainingEntries <= 0) {
        finalizeRehash(); return;
    }
    int t1Available = Math.max(1, t1.capacity() - t1.size());
    int tax = Math.max(1, (int) Math.ceil((double) t0Remaining / t1Available));

    long t0Address = t0.rawAddress();
    for (int step = 0; step < tax; step++) {
        long bAddress = t0Address + ((long) rehashCursor << 4);
        long word1 = UNSAFE.getLong(bAddress + 8);
        short distance = (short) word1;
        if (distance != EMPTY_DISTANCE) {
            long word0 = UNSAFE.getLong(bAddress);
            long origHash = t0.reconstructHash(rehashCursor, distance, 
                               (short)(word0>>>48), (int)(word1>>>32));
            t1.putByHash(origHash, (int)word0, (short)(word0>>>32));
            t0RemainingEntries--;
        }
        rehashCursor++;
    }
}
    \end{lstlisting}
    
    \vspace{1mm}
    \fontsize{8.5}{10.5}\selectfont
    \noindent\textcolor{shidenaccent}{\textbf{Rehashing Guarantee:}}\\
    \textcolor{shidenpurple}{$\bullet$ \textbf{100\% Bit Precision:}} Exact 64-bit hash reconstruction\\
    \textcolor{shidenpurple}{$\bullet$ \textbf{Multi-Engine Support:}} Supports FFM, Unsafe, \& Sharded
\end{minipage}
\clearpage

% ==============================================================================
% PAGE 10 (P8 & P9): PERFORMANCE RESULTS TABLE & CONCLUSION
% ==============================================================================
\noindent
\begin{minipage}[t]{0.48\textwidth}
    \vspace{0pt}
    {\fontsize{9}{11}\selectfont\color{shidensuccess}\textbf{P8 \& P9 // PERFORMANCE SCORECARD \& CONCLUSION}\par}
    \vspace{1mm}
    
    {\fontsize{17}{20}\selectfont\color{white}\textbf{Results \& Graduation Verdict}\par}
    \vspace{1mm}
    \noindent\rule{\linewidth}{0.4pt}
    \vspace{1.5mm}
    
    \small\color{shidentext}
    \textbf{Incremental Rehash Benchmark Metrics:}
    \begin{center}
    \fontsize{7.5}{9.5}\selectfont
    \color{shidentext}
    \begin{tabular}{lrr}
    \toprule
    \textbf{Operation / Metric} & \textbf{Baseline (No Rehash)} & \textbf{Active Rehash} \\
    \midrule
    \textbf{GET Latency} & $3.21\text{ ns}$ & $8.45\text{ ns}$ \\
    \textbf{PUT Latency} & $14.10\text{ ns}$ & $31.80\text{ ns}$ \\
    \textbf{DELETE Latency} & $12.50\text{ ns}$ & $28.40\text{ ns}$ \\
    \textbf{Tail Latency (p99.9)} & $1.10\,\mu\text{s}$ & \textbf{1.56} $\boldsymbol{\mu\text{s}}$ (vs $41.66\,\mu\text{s}$) \\
    \textbf{Hash Reconstruction} & 100\% Precision & 300,000 / 300,000 Keys \\
    \bottomrule
    \end{tabular}
    \end{center}
    
    \vspace{1mm}
    \textbf{Multi-Engine Verification Proof:}
    All 3 engine types (\texttt{SAFE\_FFM\_MONOLITHIC}, \texttt{UNSAFE\_FASTPATH\_MONOLITHIC}, and \texttt{UNSAFE\_CONTIGUOUS\_SHARDED}) successfully passed 300,000 key rehash tests with \textbf{100\% precision and 0 missing entries}.
    
    \takeaway{Dynamic incremental rehashing keeps p99.9 tail spikes under 1.56\,$\mu$s while supporting all three engine implementations with 100\% mathematical correctness.}
\end{minipage}%
\hfill
\begin{minipage}[t]{0.01\textwidth}
    \centering\color{shidenpurple!20}\rule{0.8pt}{2.1in}
\end{minipage}%
\hfill
\begin{minipage}[t]{0.48\textwidth}
    \vspace{0pt}
    {\fontsize{9}{11}\selectfont\color{shidenmuted}\textbf{FINAL ENGINEERING CONCLUSION}\par}
    \vspace{1.5mm}
    
    \fontsize{8.5}{10.5}\selectfont\color{shidentext}
    \textbf{Graduation Verdict: APPROVED FOR SHIDEN KERNEL INTEGRATION}
    
    PoC-03 successfully proves that Shiden's Partition Key Resolution Engine achieves ultra-low latency while completely bypassing JVM Garbage Collection.
    
    \vspace{2mm}
    \noindent\textcolor{shidenaccent}{\textbf{Key Architectural Accomplishments:}}
    \begin{enumerate}\setlength{\itemsep}{1pt}
      \item \textcolor{shidenpurple}{\textbf{Zero GC Overhead:}} Storing 16-byte packed buckets in off-heap native memory eliminates heap object headers and STW GC pauses entirely (\textbf{0 Bytes Allocated}).
      \item \textcolor{shidenpurple}{\textbf{Sub-10ns Lookup Speed:}} Unsafe fast-path dereferencing delivers \textbf{3.21 ns} single-key lookups and \textbf{9.56 ns} 32-key batch lookups.
      \item \textcolor{shidenpurple}{\textbf{L2 Cache Wall Resolution:}} Contiguous 512KB sharding confines probing to L2 cache, delivering a \textbf{4.8x speedup (23.34 ns)} at 1,000,000 key scale.
      \item \textcolor{shidenpurple}{\textbf{Flat Tail Latency:}} Adaptive incremental rehashing limits p99.9 latency spikes to \textbf{1.56 $\mu$s} across all three engine backends.
    \end{enumerate}
    
    \vspace{2mm}
    \begin{center}
      \fontsize{9}{11}\selectfont\color{shidensuccess}
      \textbf{STATUS: PoC-03 READY FOR RFC-002 PRODUCTION COMPOSITION}
    \end{center}
\end{minipage}

\end{document}
